\documentclass[letterpaper,11pt]{article}

\usepackage{latexsym}
\usepackage[empty]{fullpage}
\usepackage{titlesec}
\usepackage{marvosym}
\usepackage[usenames,dvipsnames]{color}
\usepackage{verbatim}
\usepackage{enumitem}
\usepackage[hidelinks]{hyperref}
\usepackage{fancyhdr}
\usepackage[english]{babel}
\usepackage{tabularx}
\input{glyphtounicode}

\pagestyle{fancy}
\fancyhf{}
\fancyfoot{}
\renewcommand{\headrulewidth}{0pt}
\renewcommand{\footrulewidth}{0pt}

\addtolength{\oddsidemargin}{-0.5in}
\addtolength{\evensidemargin}{-0.5in}
\addtolength{\textwidth}{1in}
\addtolength{\topmargin}{-.5in}
\addtolength{\textheight}{1.0in}

\urlstyle{same}
\raggedbottom
\raggedright
\setlength{\tabcolsep}{0in}

\titleformat{\section}{
  \vspace{-4pt}\scshape\raggedright\large
}{}{0em}{}[\color{black}\titlerule \vspace{-5pt}]

\pdfgentounicode=1

\newcommand{\resumeItem}[1]{
  \item\small{
    {#1 \vspace{-4pt}}
  }
}
\newcommand{\resumeSubheading}[4]{
  \vspace{-2pt}\item
    \begin{tabular*}{0.97\textwidth}[t]{l@{\extracolsep{\fill}}r}
      \textbf{#1} & #2 \\
      \textit{\small#3} & \textit{\small #4} \\
    \end{tabular*}\vspace{-7pt}
}
\newcommand{\resumeSubSubheading}[2]{
    \item
    \begin{tabular*}{0.97\textwidth}{l@{\extracolsep{\fill}}r}
      \textit{\small#1} & \textit{\small #2} \\
    \end{tabular*}\vspace{-7pt}
}
\newcommand{\resumeProjectHeading}[2]{
    \item
    \begin{tabular*}{0.97\textwidth}{l@{\extracolsep{\fill}}r}
      \small#1 & #2 \\
    \end{tabular*}\vspace{-7pt}
}
\newcommand{\resumeSubItem}[1]{\resumeItem{#1}\vspace{-4pt}}
\renewcommand\labelitemii{$\vcenter{\hbox{\tiny$\bullet$}}$}
\newcommand{\resumeSubHeadingListStart}{\begin{itemize}[leftmargin=0.15in, label={}]}
\newcommand{\resumeSubHeadingListEnd}{\end{itemize}}
\newcommand{\resumeItemListStart}{\begin{itemize}}
\newcommand{\resumeItemListEnd}{\end{itemize}\vspace{-5pt}}

%-------------------------------------------

\begin{document}

\begin{center}
    \textbf{\Huge \scshape Yali Sommer} \\ \vspace{1pt}
    \small 954-901-9375 $|$
    \href{mailto:yali_sommer@brown.edu}{\underline{yali\_sommer@brown.edu}} $|$
    \href{https://linkedin.com/in/yalisommer}{\underline{LinkedIn}} $|$
    \href{https://github.com/yalisommer}{\underline{GitHub}} $|$
    \href{https://yalisommer.com}{\underline{yalisommer.com}} \\[3pt]
    \small \textbf{Computer Vision $|$ Differentiable Rendering $|$ Geometry Processing} $|$
    C++, OpenGL, GLSL, PyTorch, Mitsuba 3, BlenderPy
\end{center}

%-----------EDUCATION-----------
\section{Education}
  \resumeSubHeadingListStart
    \resumeSubheading
      {Brown University $|$ 3.94 GPA}{Providence, RI}
      {Computer Science Sc.B., Mathematics A.B.}{Sept. 2023 -- May 2027}
      \resumeSubHeadingListStart
        \resumeItem{\textbf{Coursework:} Computer Vision, Computer Graphics, Computer Graphics: Geometry \& Simulation, Data Structures \& Algorithms, Computer Systems, Abstract Algebra, Complex Analysis}
        \resumeItem{\textbf{Teaching:} Undergraduate TA --- Introduction to Computer Graphics (CSCI 1230); Data Structures \& Algorithms (CSCI 0200)}
      \resumeSubHeadingListEnd
  \resumeSubHeadingListEnd

%-----------RESEARCH-----------
\section{Research}
\resumeSubHeadingListStart

  \resumeSubheading
    {Brown Visual Computing Lab}{Aug. 2025 -- Present}
    {Research Assistant $|$ Supervisors: Prof. James Tompkin \& Joel Salzman (PhD)}{Providence, RI}
    \resumeItemListStart
      \resumeItem{Recovering curved mirror geometry from catacaustic light patterns via gradient-based optimization over a differentiable Mitsuba 3 rendering pipeline, treating mirror coefficients as free parameters fit to observed reflected-light distributions.}
      \resumeItem{PyTorch optimizer with MS-SSIM perceptual loss and finite-difference gradient estimation to fit parabolic mirror coefficients to rendered scene observations.}
      \resumeItem{Validated across planar and parabolic configurations using synthetic BlenderPy scenes.}
    \resumeItemListEnd

  \resumeSubheading
    {University of Edinburgh --- Geometry Processing Group}{2025 -- Present}
    {Research Collaborator $|$ Supervisor: Prof. Amir Vaxman}{Edinburgh, UK $|$ Remote}
    \resumeItemListStart
      \resumeItem{Adding constraint satisfaction to the DROK mesh deformation framework: quad planarity and edge stretch conditions encoded directly into training rather than applied as post-corrections.}
      \resumeItem{Lagrangian multiplier methods and penalty training enforce these constraints during learning, preserving deformation quality across mesh topologies without a correction pass.}
      \resumeItem{Built a real-time constraint-aware mesh editor within the DROK deformation space; continuing as Brown Senior Thesis co-advised by Prof. Adriana Schulz and Prof. Daniel Ritchie.}
    \resumeItemListEnd

\resumeSubHeadingListEnd

%-----------PROJECTS-----------
\section{Projects}
\resumeSubHeadingListStart

  \resumeProjectHeading
    {\textbf{Realtime OpenGL Renderer} $|$ \emph{C++, OpenGL, GLSL}}{Oct. 2025 -- Dec. 2025}
    \resumeItemListStart
      \resumeItem{Built a real-time C++/OpenGL renderer with a modular GLSL shader pipeline, scene management, and camera controls.}
      \resumeItem{Implemented percentage-closer soft shadow mapping, bump mapping, and a GPU compute shader particle system supporting fire, smoke, and snow effects.}
      \resumeItem{Added a full-screen post-processing pass with depth-of-field, fog, and stylized line-art filters driven by Z-buffer readback.}
      \resumeItem{Maintained 60+ FPS at high polygon counts through draw-call batching and GPU instancing.}
    \resumeItemListEnd

  \resumeProjectHeading
    {\textbf{Ray Tracer} $|$ \emph{C++, Ray Tracing, BVH}}{Jan. 2025 -- Mar. 2025}
    \resumeItemListStart
      \resumeItem{Built a CPU ray tracer in C++ with Phong shading, area-light soft shadows, multi-bounce reflections, and Snell's law refractions.}
      \resumeItem{Added stochastic anti-aliasing and depth of field via aperture sampling.}
      \resumeItem{Implemented a BVH acceleration structure from scratch to reduce intersection cost across triangle mesh scenes.}
    \resumeItemListEnd

  \resumeProjectHeading
    {\textbf{In-Browser Fish Detection} $|$ \emph{TypeScript, React, YOLOv8, ONNX Runtime Web}}{2025 -- Present}
    \resumeItemListStart
      \resumeItem{Trained a YOLOv8 fish species detector from scratch on a HuggingFace dataset and deployed it entirely client-side via ONNX Runtime Web, running at $\sim$10 FPS as a WebAssembly artifact.}
      \resumeItem{Built the full TypeScript inference stack: GPU-threaded ONNX sessions, frame preprocessing, NMS post-processing, and canvas overlay rendering.}
    \resumeItemListEnd

  \resumeProjectHeading
    {\textbf{ViolenceNet} $|$ \emph{Python, PyTorch, TensorFlow, OpenCV, YOLO}}{Mar. 2025 -- May 2025}
    \resumeItemListStart
      \resumeItem{Designed a custom 3D CNN on stacked RGB frame volumes for violence classification (94\% val accuracy); integrated fine-tuned YOLOv8 for concurrent dangerous-item detection.}
      \resumeItem{Evaluated against videos flagged by YouTube's content moderation system; 88\% accuracy on real-world moderated content.}
    \resumeItemListEnd

\resumeSubHeadingListEnd

%-----------EXPERIENCE-----------
\section{Experience}
\resumeSubHeadingListStart

  \resumeSubheading
    {IturanTech}{May 2025 -- Aug. 2025}
    {AI/ML Intern}{Tel Aviv, Israel}
    \resumeItemListStart
      \resumeItem{Engineered features for an LSTM anomaly detection pipeline across 1M+ connected vehicle units; built a regex aggregation layer translating raw telematics readings to structured text narratives.}
      \resumeItem{Fine-tuned an LLM on labeled fleet anomaly data for the same detection task; presented to VP of Engineering and advanced for post-internship development.}
    \resumeItemListEnd

  \resumeSubheading
    {SMBC}{June 2026 -- Aug. 2026}
    {AI Intern, Data Strategy Team}{New York, NY}
    \resumeItemListStart
      \resumeItem{Incoming AI Intern on SMBC's Data Strategy team.}
    \resumeItemListEnd

\resumeSubHeadingListEnd

%-----------SKILLS-----------
\section{Skills}
  \resumeSubHeadingListStart
    \resumeItem{\textbf{Graphics \& Vision:} OpenGL, GLSL, Mitsuba 3, BlenderPy, OpenCV, ONNX Runtime Web} \vspace{-1mm}
    \resumeItem{\textbf{Languages:} C++, Python, GLSL, TypeScript, SQL, Java} \vspace{-1mm}
    \resumeItem{\textbf{ML \& Research:} PyTorch, WandB, NumPy, Differentiable Rendering, Gradient-Based Optimization, Inverse Methods} \vspace{-1mm}
  \resumeSubHeadingListEnd

\end{document}
